% !TEX program = xelatex
% Copyright (c) 2026 Ingolf Lohmann.
% External publication or relicensing requires a separate decision.

\documentclass[11pt,a4paper]{article}

\usepackage[a4paper,margin=23mm,headheight=15pt]{geometry}
\usepackage[provide=*,ngerman]{babel}
\usepackage{fontspec}
\usepackage{microtype}
\usepackage{amsmath,amssymb,mathtools}
\usepackage{booktabs,longtable,tabularx,array}
\usepackage{enumitem}
\usepackage{xcolor}
\usepackage{listings}
\usepackage{seqsplit}
\usepackage{fancyhdr}
\usepackage{hyperref}

\setmonofont{DejaVu Sans Mono}

\definecolor{qikblue}{HTML}{123A63}
\definecolor{qikgold}{HTML}{A26A00}
\definecolor{qikgreen}{HTML}{176B3A}
\definecolor{qikred}{HTML}{9D2929}
\definecolor{qikgray}{HTML}{F2F4F7}
\definecolor{qikdarkgray}{HTML}{4A5561}

\hypersetup{
  colorlinks=true,
  linkcolor=qikblue,
  urlcolor=qikblue,
  citecolor=qikgreen,
  pdftitle={Von der Planck-Bruecke zur massiven Kugel -- QIK-VRT SMG H5},
  pdfauthor={Ingolf Lohmann},
  pdfsubject={Kernelgepruefter Modellkern, offene Physikbruecken und virtuelle Kosmogenese},
  pdfkeywords={QIK-VRT, Planck, Standardmodell, Gravitation, Lean, Kausalitaet, virtuelle Kosmogenese}
}

\pagestyle{fancy}
\fancyhf{}
\lhead{QIK--VRT · SMG H5}
\rhead{Ingolf Lohmann · 2. August 2026}
\cfoot{\thepage}

\setlist[itemize,enumerate]{nosep,leftmargin=*}
\setlength{\parindent}{0pt}
\setlength{\parskip}{0.54em}
\setlength{\emergencystretch}{3em}

\newcommand{\QIKVRT}{QIK--VRT}
\newcommand{\SMGVRT}{\ensuremath{\mathrm{SMG}_{\mathrm{VRT}}}}
\newcommand{\code}[1]{\texttt{#1}}
\newcommand{\openstatus}{\textcolor{qikred}{\textbf{OPEN}}}
\newcommand{\kernelstatus}{\textcolor{qikgreen}{\textbf{KERNEL-GEPRÜFT}}}
\newcommand{\conditionalstatus}{\textcolor{qikgold}{\textbf{BEDINGT}}}
\newcommand{\leanname}[1]{{\ttfamily\footnotesize\seqsplit{#1}}}

\newenvironment{statusbox}{%
  \begin{center}\begin{minipage}{0.94\textwidth}\small
  \color{qikdarkgray}\hrule\vspace{0.7em}
}{%
  \vspace{0.7em}\hrule\end{minipage}\end{center}
}

\newenvironment{claimbox}{%
  \begin{center}\begin{minipage}{0.92\textwidth}
  \color{qikblue}\hrule\vspace{0.7em}\color{black}
}{%
  \vspace{0.7em}\color{qikblue}\hrule\end{minipage}\end{center}
}

\lstdefinelanguage{Lean}{
  morekeywords={import,namespace,end,inductive,structure,where,def,theorem,by,
    cases,decide,rfl,simp,simpa,using,exact,fun,Type,Prop,Bool,Nat,Int,
    true,false,none,some,deriving,if,then,else},
  sensitive=true,
  morecomment=[l]{--},
  morecomment=[s]{/-}{-/},
  morestring=[b]"
}

\lstset{
  basicstyle=\ttfamily\footnotesize,
  keywordstyle=\color{qikblue}\bfseries,
  commentstyle=\color{qikdarkgray}\itshape,
  stringstyle=\color{qikgreen},
  backgroundcolor=\color{qikgray},
  frame=single,
  rulecolor=\color{qikdarkgray},
  breaklines=true,
  columns=fullflexible,
  keepspaces=true,
  showstringspaces=false,
  tabsize=2
}

\title{\textbf{Von der Planck-Brücke zur massiven Kugel}\\[0.4em]
\large \QIKVRT{} SMG H5: kernelgeprüfter Modellkern,\\
\large offene Physikbrücken und virtuelle Kosmogenese}
\author{Ingolf Lohmann}
\date{2. August 2026}

\begin{document}

\maketitle

\begin{statusbox}
\textbf{Dokumentstatus:} Fachartikel- und Peer-Review-Kandidat; nicht extern publiziert.\\
\textbf{Formaler Status:} 32 von 32 H5-Theoremen mit Lean 4.19.0 kernelakzeptiert;
keine projektspezifischen Axiome, kein \code{sorry}, \code{admit} oder
\code{unsafe}.\\
\textbf{Axiom-Audit:} 17 Theoreme ohne Axiome; 13 nur mit \code{propext};
2 arithmetische Wachstumssätze mit \code{propext} und \code{Quot.sound}.\\
\textbf{Physikalischer Status:} \code{PHYSICAL\_UNIFICATION = OPEN\_CANDIDATE};
keine Behauptung eines beobachteten Gravitons, einer fertigen Quantengravitation
oder eines physisch erzeugten Urknalls.\\
\textbf{Wirkungsstatus:} \code{EFFECT\_ACK\_CONTINUE}; kein repositoryweites
\code{PASS}, \code{FINAL\_PASS} oder \code{EFFECT\_ACK\_DONE}.
\end{statusbox}

\begin{abstract}
\QIKVRT{} SMG H5 formuliert einen additiven, fail-closed Prüfrahmen für den
Versuch, relationale Kausalität an Standardmodell, Allgemeine
Relativitätstheorie, Planck-Skala und Quantenfeldsprache anzuschließen. Der
Lean-Kern prüft die symbolische Planck-Normalform, die gemeinsame
Ereignisidentität von Wave- und Record-Projektionen, getrennte Evidenzgates,
eine zwölfteilige Schließungsbedingung sowie monotone und bedingt
unbeschränkte Expansion eines virtuellen Übergangssystems. Derselbe Kern weist
den aktuellen physikalischen Kandidaten als nicht massiv geschlossen aus.
Standardmodell- und Einstein-Grenzfall, universelle Kopplung einschließlich
Higgs-Sektor, Quantengravitationskorrespondenz, Stabilität, Unitarität,
Nichtzirkularität, unterscheidende Vorhersage, empirische Korrespondenz und
unabhängige Reproduktion bleiben offen. H5 ist damit ein substantieller
maschinenprüfbarer Brücken- und Falsifikationskandidat, keine abgeschlossene
Entdeckung einer neuen Naturtheorie.
\end{abstract}

\tableofcontents

\section{Leitthese und Kreisstruktur}

\begin{claimbox}
\begin{center}
\Large\textbf{Kausalität ist Relation, nicht Sequenz.}
\end{center}
\end{claimbox}

Eine zeitliche Reihenfolge kann Teil einer Kausalbeschreibung sein, ist aber
kein hinreichender Kausalbeweis. Erst eine nachprüfbare Abhängigkeit zwischen
Quelle, Transformation, Kontext, Entscheidung, Autorisierung und Wirkung macht
aus einem bloßen ``danach'' ein begründetes ``deshalb''.

Die QIK-VRT-Kette lautet:

\[
\begin{aligned}
\text{Unterscheidung}
&\to\text{Information}
\to\text{Messung}
\to\text{Wirkung}
\to\text{Relation}\\
&\to\text{Kausalordnung}
\to\text{rekonstruierbare Raumzeit}
\to\text{klassischer Grenzfall}.
\end{aligned}
\]

Der Kreis schließt sich, weil in der emergierten Welt Menschen und Maschinen
wieder Unterschiede erzeugen, messen, speichern und in Wirkung überführen.
Damit kehrt die physikalische Frage zur Verantwortungsfrage zurück: Welche
Relationen rechtfertigen eine Wirkung, und welche halten sie an?

\subsection{Epistemische Trennung}

\begin{longtable}{>{\bfseries}p{0.22\textwidth}p{0.34\textwidth}p{0.34\textwidth}}
\toprule
Klasse & Zulässige Aussage & Unzulässige Aufwertung \\
\midrule
\endhead
formal bewiesen & Ein Kernel prüft Folgen expliziter Definitionen. & Die Natur erfüllt deshalb das Modell. \\
empirisch gestützt & Ein Messbefund trägt im publizierten Scope. & Er beweist jede weitergehende Ontologie. \\
quellengebunden & Eine Quelle stützt eine begrenzte Aussage. & Autorität ersetzt Herleitung oder Messung. \\
normativ & Eine Regel bindet Verantwortung und Wirkung. & Die Regel folgt zwingend aus Physik. \\
interpretativ & Ein Rahmen ordnet Befunde und Begriffe. & Die Deutung ist ein Naturtheorem. \\
offen & Eine Brücke besitzt Frage und Schließungsbedingung. & Offenheit wird als stiller PASS behandelt. \\
\bottomrule
\end{longtable}

\section{Relation, Kausalstruktur und Raumzeit}

Die relationale Ausgangsthese ist physikalisch anschlussfähig, ohne durch
bestehende Arbeiten bereits bewiesen zu sein. Hawking, King und McCarthy sowie
Malament zeigten unter ausdrücklichen Raumzeitvoraussetzungen, wie viel
topologische und konforme Struktur aus Kausal- beziehungsweise zeitartigen
Relationen rekonstruierbar ist \cite{hkm1976,malament1977}. Causal-Set-Modelle
untersuchen lokal endliche partielle Ordnungen als mögliche Fundamentalstruktur
\cite{bombelli1987}.

Der Prozessmatrix-Formalismus erlaubt lokale Quantenoperationen ohne eine
vorausgesetzte einzige globale Kausalordnung \cite{oreshkov2012}. Kohärent
kontrollierte Gate-Reihenfolgen wurden experimentell demonstriert
\cite{procopio2015}. Daraus folgt keine frei wählbare Nachricht in einen bereits
persistierten früheren Record. Es folgt die methodische Vorsicht, globale
klassische Sequenz nicht unnötig früh als Fundamentalstruktur zu setzen.

\section{Wave und Record}

H5 liest die Wave-Seite als Möglichkeits-, Ausbreitungs- und
Amplitudenstruktur und die Record-Seite als lokalisierten Wechselwirkungsausgang.
Ein typisiertes Objekt

\[
D=(\operatorname{id},W,R)
\]

besitzt die Projektionen

\[
\pi_W(D)=(\operatorname{id},W),
\qquad
\pi_R(D)=(\operatorname{id},R).
\]

Der Kernel beweist:

\[
\operatorname{id}(\pi_W(D))
=\operatorname{id}(D)
=\operatorname{id}(\pi_R(D)).
\]

Dies ist eine Repräsentationsinvariante, keine Festlegung der einzig möglichen
Quantenontologie.

\subsection{Higgs als empirischer Anker}

In der Quantenfeldsprache ist das Higgs-Boson eine Anregung des Higgs-Feldes
\cite{cernhiggs}. ATLAS und CMS beobachteten 2012 unabhängig einen neuen
Boson-Zustand um 125 GeV \cite{atlas2012,cms2012}. Dies stützt das allgemeine
Muster Feld--Anregung--lokalisierter Detektorrecord. Es beobachtet weder ein
Graviton noch leitet es eine konkrete Gravitationstheorie her.

\section{Gravitation: gemessene Welle, offenes Quant}

Mit GW150914 wurde eine Gravitationswelle direkt beobachtet; die gemessene
Wellenform stimmte im geprüften Bereich mit der Allgemeinen Relativitätstheorie
überein \cite{ligo2016}. Ein einzelnes Graviton ist nicht direkt beobachtet.
Das Standardmodell enthält die Gravitation nicht \cite{cernsm}.

H5 bindet den Ausgangsstand:

\begin{lstlisting}
higgsFieldExcitationObserved = true
gravitationalWaveObserved = true
gravitonObserved = false
quantumGravityPredictionConfirmed = false
\end{lstlisting}

Die Theoreme H5-T12 bis H5-T14 verhindern die automatische Evidenzhochsetzung.
Sie beweisen nicht, dass Gravitonen unmöglich sind.

\section{Planck-Normalform}

Für reduzierte Compton-Wellenlänge und Gravitationsradius gilt bei der
Planck-Masse:

\[
\bar\lambda_C(m_P)
=\frac{\hbar}{m_Pc}
=\ell_P
=\frac{Gm_P}{c^2}
=r_g(m_P).
\]

Der Schwarzschild-Radius ist davon durch den Faktor zwei zu unterscheiden:

\[
r_S(m_P)=\frac{2Gm_P}{c^2}=2\ell_P.
\]

Weiter gelten symbolisch

\[
\ell_Pp_P=t_PE_P=\hbar,
\qquad
\frac{\ell_P}{t_P}=\frac{E_P}{p_P}=c.
\]

\subsection{Exakte Kernelrepräsentation}

H5 repräsentiert Monome in \((\hbar,G,c)\) durch verdoppelte ganzzahlige
Exponenten. Das Tripel \((1,1,-3)\) steht für

\[
\hbar^{1/2}G^{1/2}c^{-3/2}=\ell_P.
\]

Die Basismonome lauten:

\[
\ell_P=(1,1,-3),\quad
t_P=(1,1,-5),\quad
m_P=(1,-1,1).
\]

Multiplikation addiert, Division subtrahiert die Tripel. So prüft Lean die
Normalform exakt, ohne Gleitkommarundung.

\begin{longtable}{p{0.13\textwidth}p{0.57\textwidth}p{0.2\textwidth}}
\toprule
Satz & Kernelgeprüfte Aussage & Axiome \\
\midrule
\endhead
H5-T01 & \(\hbar/(m_Pc)=\ell_P\) als Monomidentität & keine \\
H5-T02 & \(Gm_P/c^2=\ell_P\) als Monomidentität & keine \\
H5-T03 & \(\ell_Pp_P=\hbar\) als Monomidentität & keine \\
H5-T04 & \(t_PE_P=\hbar\) als Monomidentität & keine \\
H5-T05 & \(\ell_P/t_P=c\) als Monomidentität & keine \\
H5-T06 & \(E_P/p_P=c\) als Monomidentität & keine \\
H5-T07 & gemeinsame sechsteilige Normalform & keine \\
\bottomrule
\end{longtable}

Die Identitäten markieren eine natürliche Skalenbeziehung. Sie liefern noch
keine mikroskopische Dynamik, keinen physikalischen Zustandsraum und keine
bestätigte UV-Vervollständigung.

\section{SMG-VRT: Erweiterung des Standardmodells}

Weil Gravitation nicht Teil des Standardmodells ist, bezeichnet H5 den
Arbeitsrahmen als Erweiterung \SMGVRT. Ein schematischer Anschluss lautet:

\[
S_{\SMGVRT}
=S_{\mathrm{SM}}[g,\Phi]
+\frac{c^3}{16\pi G}\int d^4x\sqrt{-g}\,R
+S_{\mathrm{bridge}}[g,\Phi,\chi].
\]

Hier sind \(\Phi\) die Standardmodellfelder einschließlich Higgs-Feld,
\(g\) die gravitative oder rekonstruierte geometrische Struktur und \(\chi\)
zusätzliche relationale Freiheitsgrade. Der bekannte Standardmodell- plus
Einstein-Hilbert-Anteil ist noch keine neue Vereinheitlichung. Allgemeine
Relativität besitzt eine konsistente Niederenergie-Effektivfeldtheorie
\cite{donoghue1994}; die Hochenergievervollständigung bleibt davon
unterbestimmt. Die Konsistenzroute von einem linearen masselosen
Spin-2-Gaugefeld zur Einstein-Selbstkopplung ist ebenfalls bekannt
\cite{deser1970}.

Der neue physikalische Gehalt müsste in einem expliziten
\(S_{\mathrm{bridge}}\), seinen Symmetrien, Observablen und unterscheidenden
Vorhersagen liegen.

\subsection{Erforderliche Grenzsätze}

Ein geschlossener Kandidat muss mindestens zeigen:

\[
\SMGVRT
\xrightarrow[\text{niedrige gravitative Kopplung}]{}
\mathrm{SM}
\]

im geprüften Teilchen-, Symmetrie- und Amplitudenscope sowie

\[
\SMGVRT
\xrightarrow[\text{klassisch / grobgekörnt}]{}
G_{\mu\nu}+\Lambda g_{\mu\nu}
=\frac{8\pi G}{c^4}T_{\mu\nu}.
\]

Die universelle Kopplung muss den Energie-Impuls-Beitrag des Higgs-Sektors
einschließen. Semantische Information darf nicht ohne physischen Träger als
zusätzliche Energiequelle eingesetzt werden.

\section{Zwölfteilige massive Schließung}

H5 definiert:

\[
\begin{aligned}
\operatorname{MassiveClosure}(M):={}&P\land D\land S\land E\land U\land Q\\
&\land V\land C\land N\land F\land X\land R.
\end{aligned}
\]

\begin{longtable}{>{\bfseries}p{0.08\textwidth}p{0.58\textwidth}p{0.23\textwidth}}
\toprule
Code & Zeuge & aktueller H5-Stand \\
\midrule
\endhead
P & Planck-Normalform & \kernelstatus{} im Symbolmodell \\
D & Wave/Record-Identität & \kernelstatus{} im Typmodell \\
S & Standardmodell-Grenzfall & \openstatus \\
E & klassischer Einstein-Grenzfall & \openstatus \\
U & universelle Stress-Energie-Kopplung inkl. Higgs & \openstatus \\
Q & Quantengravitationskorrespondenz & \openstatus \\
V & Stabilität und Unitarität & \openstatus \\
C & kausale Konsistenz & \kernelstatus{} im H5-Scope \\
N & Nichtzirkularität & \openstatus \\
F & falsifizierbare unterscheidende Vorhersage & \openstatus \\
X & empirische Korrespondenz & \openstatus \\
R & unabhängige Reproduktion & \openstatus \\
\bottomrule
\end{longtable}

Die Theoreme H5-T15 bis H5-T24 zeigen, dass Schließung die kritischen Zeugen
voraussetzt. H5-T25 bestätigt:

\begin{lstlisting}
massiveClosure currentH5Candidate = false
\end{lstlisting}

H5-T26 zeigt lediglich die konstruktive Erfüllbarkeit der Booleschen
Schließungsfunktion durch vollständig gelieferte Daten. Er liefert diese
physischen Daten nicht.

\section{Kernel-Receipt ist keine Naturentdeckung}

Der technische Receipt besitzt vier Felder: Quellbindung, Elaboration,
Kernelannahme und Axiom-Audit. Selbst wenn alle vier wahr sind, bleibt der
Status ohne physikalische Schließung auf \code{formalModelChecked}.

H5-T27 beweist diesen konkreten Status. H5-T28 beweist umgekehrt, dass die im
Modell so bezeichnete empirische Korroboration sowohl vollständigen
Kernel-Receipt als auch massive Schließung voraussetzt.

\begin{claimbox}
\textbf{Grenze:} Ein formaler Kernel prüft Mathematik relativ zu Definitionen.
Er ist kein Detektor und kann die Instanziierung seiner Variablen durch die
Natur nicht selbst feststellen.
\end{claimbox}

\section{Virtuelle Kosmogenese}

``Urknall in einer virtuellen Maschine'' wird operational als Seed plus
typisierte Übergangsrelation definiert. Für Seed \(S_0\) und
\(T\subseteq S\times S\) gilt:

\[
\begin{aligned}
R_0(x)&:\Leftrightarrow x=S_0,\\
R_{n+1}(x)&:\Leftrightarrow R_n(x)
\lor\exists y\,[R_n(y)\land T(y,x)].
\end{aligned}
\]

Lean beweist die Monotonie

\[
R_n(x)\Rightarrow R_{n+1}(x)
\]

und die Erreichbarkeit des Seeds in jeder endlichen Stufe. Für eine
Populationsfunktion \(P:\mathbb N\to\mathbb N\) und einen gelieferten
strikten Wachstumszeugen gilt außerdem:

\[
\Bigl[\exists n_0\;\forall n\ge n_0:\;P(n)<P(n+1)\Bigr]
\Rightarrow
\Bigl[\forall B\;\exists n:\;B<P(n)\Bigr].
\]

Dies ist ein Theorem über ein virtuelles Modell. Es beweist weder die Identität
mit dem physischen frühen Universum noch reale Materie-, Raumzeit- oder
Energieerzeugung durch einen Computer.

\section{Syntax, Semantik und Prüfkette}

Die EBNF-Oberfläche enthält Pflichtblöcke für Planck-Brücke, Dualität,
empirische Anker, Grenzfälle, Kopplung, Konsistenz, Vorhersage, Schließung,
virtuelle Kosmogenese, Effektgrenze und offene Obligationen.

Der Referenzvalidator prüft zwanzig statische Regeln. Elf positive und negative
Tests blockieren unter anderem:

\begin{itemize}
  \item fehlende oder doppelte Pflichtblöcke;
  \item die Verwechslung von Gravitations- und Schwarzschild-Radius;
  \item Wave/Record-Identität als vermeintliche physische Vollständigkeit;
  \item Gravitationswelle als vermeintlichen Gravitonnachweis;
  \item massive Schließung bei mindestens einem falschen Zeugen;
  \item Kernelannahme als physikalische Entdeckung;
  \item virtuelle Kosmogenese als physischer Urknall;
  \item \code{EFFECT\_ACK\_DONE} aus dem wissenschaftlichen Kandidaten.
\end{itemize}

\begin{lstlisting}
EBNF source
-> parsed typed blocks
-> static obligations S01-S20
-> Lean model semantics
-> kernel receipt and axiom audit
-> claim/source matrix
-> EFFECT_ACK_CONTINUE
\end{lstlisting}

Die separate Vektorgrafik
\code{VRTCore\_SMG\_EBNF\_Map\_DE.svg} bildet diese Kreisstruktur ab.

\section{Offenes Dynamik- und Falsifikationsprogramm}

Die physikalische Lücke wird durch folgende Arbeitskette geschlossen oder
widerlegt:

\begin{enumerate}
  \item vollständige Feld-, Zustands- und Observabledefinition;
  \item lokale Wirkung \(S_{\mathrm{bridge}}\);
  \item Symmetrie-, Zwangs- und Anomalieanalyse;
  \item Standardmodell-Decoupling-Theorem;
  \item Einstein-/Newton-Grenztheorem;
  \item Higgs-inclusive universelle Kopplung;
  \item linearisierter Wellensektor und definierter Quantenzustand;
  \item Ghost-, Stabilitäts-, Unitaritäts- und Kausalitätsanalyse;
  \item nichtzirkuläre Abbildung auf Messoperationen;
  \item eingefrorene numerische Vorhersage;
  \item unabhängiger Rechen- und Datenbenchmark;
  \item erst danach physikalischer Discovery-Review.
\end{enumerate}

Eine unterscheidende Vorhersage muss mindestens die Form

\[
\Delta O
=O_{\SMGVRT}-O_{\mathrm{SM+GR/EFT}}
\neq0
\]

für ein zugängliches Observable \(O\) besitzen. Vor Messung werden Parameter,
Unsicherheit, Datenpipeline und Widerlegungskriterium eingefroren.

\section{Wissenschaftliche und technologische Tragweite}

\subsection{Erreichte Leistung}

Es ist eine große formale und konzeptionelle Leistung, einen über mehr als ein
Jahr entwickelten Gedanken von der Alltagssprache über Repository-Provenienz,
Kausalitäts- und Wirkungsgates bis zu EBNF, statischer Semantik und einem
tatsächlich ausgeführten Lean-Kern zu führen. Darauf darf der Autor stolz sein.
Die Leistung wird durch die offenen Brücken nicht entwertet; ihre klare
Kennzeichnung macht sie wissenschaftlich belastbarer.

\subsection{Bedingte Größenordnung}

Falls ein explizites \(S_{\mathrm{bridge}}\) alle zwölf Schließungszeugen
liefert, getestete Grenzfälle reproduziert und eine neue Vorhersage unabhängig
bestätigt wird, wäre die Größenordnung außerordentlich: Dann läge ein Kandidat
für eine gemeinsame Beschreibung von Quantenfeldstruktur, Gravitation,
Kausalordnung und Raumzeitrekonstruktion vor. Dies ist heute eine begründete
Möglichkeit, keine bestätigte Feldbewertung.

\subsection{Technologische Horizonte}

Unabhängig von einer späteren physikalischen Bestätigung sind anschlussfähig:

\begin{itemize}
  \item beweisstatusbewusste wissenschaftliche Repositories;
  \item KI-Systeme mit Trennung von Berechnung und Wirkung;
  \item auditierbare Quanten-klassische Kontrollschleifen;
  \item kausal adressierte digitale Zwillinge;
  \item formal gebundene Softwarelieferketten;
  \item reproduzierbare Modell- und Messkorrespondenz;
  \item Simulationen mit expliziten Emergenz- und Falsifikationsbedingungen.
\end{itemize}

\section{Mensch, Weltanschauung und Spiritualität}

\subsection{Freiheit als Anschlussfähigkeit}

Der Mensch steht nicht außerhalb der Kausalität. Freiheit muss dennoch nicht
zur Illusion werden. Sie kann als reale, verkörperte Anschlussfähigkeit
verstanden werden: unterscheiden, erinnern, bewerten, lernen, widersprechen,
korrigieren und verantwortlich handeln.

\begin{claimbox}
\textbf{Normative QIK-VRT-Linie:} Technische Möglichkeit ist noch keine
verantwortete Wirkung.
\end{claimbox}

Mit der möglichen Bedeutung einer Arbeit wächst die Verantwortung gegenüber
dem eigenen Wohlergehen, dem persönlichen Umfeld, Fachleuten, Öffentlichkeit
und Wahrheit. Unabhängige Kritik ist kein Angriff auf die Idee, sondern Teil
ihrer Wirkungsfähigkeit.

\subsection{Weltanschauliche Lesart}

QIK-VRT legt eine relationale Sicht nahe: Identität umfasst auch die
Kontinuität eines rekonstruierbaren Zusammenhangs; Wissen umfasst Quelle,
Messung, Modell, Status, Grenze und Korrekturmöglichkeit. Diese Folgerung ist
interpretativ, nicht durch Lean als Metaphysik erzwungen.

\subsection{Spirituelle Möglichkeit}

H5 beweist keine Religion, keinen Gott und keinen letzten Sinn. Eine
spirituelle Lesart kann dennoch in Verbundenheit, Staunen und Verantwortung
liegen:

\begin{quote}
Unterschiede ermöglichen Beziehung; Beziehung trägt Wirkung; Wirkung begründet
Verantwortung.
\end{quote}

Die Aussage darf existenziell bedeutsam sein, bleibt aber eine Deutung und kein
Messbefund.

\section{Ergebnis und Rückkehr zum Anfang}

H5 schließt nicht die gesamte Physik. Es schließt eine entscheidende
Vorstufe: Formalität, Empirie, physikalische Brücke, Interpretation, Norm und
Wirkung können nicht mehr still ineinander geschoben werden.

\begin{itemize}
  \item Die Planck-Normalform ist symbolisch kernelgeprüft.
  \item Wave und Record bewahren eine gemeinsame Ereignisidentität.
  \item Higgs- und Gravitationswellenbefunde sind starke Anker, aber kein
        Gravitonnachweis.
  \item Gravitation verlangt eine explizite Erweiterung des Standardmodells.
  \item Massive Schließung verlangt zwölf getrennte Zeugen.
  \item Der aktuelle physikalische Kandidat bleibt offen.
  \item Virtuelle Kosmogenese besitzt einen beweisbaren monotonen Kern, ist
        aber nicht mit dem physischen Urknall identifiziert.
\end{itemize}

Der Kreis führt vom ersten Unterschied über Universum, Physik, Mathematik,
Informatik, Mensch und Verantwortung zurück zum Unterschied zwischen möglicher
und verantworteter Wirkung.

\begin{claimbox}
Jede Behauptung und jede Wirkung soll ihre unterscheidbaren Voraussetzungen,
ihre Relationen, ihre Grenzen und ihre Verantwortung mitführen.
\end{claimbox}

\begin{center}
\Large\textbf{Kausalität ist Relation, nicht Sequenz.}
\end{center}

\section*{Maschinenlesbarer Status}

\begin{lstlisting}
FORMAL_MODEL_CORE = KERNEL_ACCEPTED_32_OF_32
PROJECT_AXIOMS = NONE
SMG_VRT_DYNAMICS = OPEN
PHYSICAL_UNIFICATION = OPEN_CANDIDATE
GRAVITON_OBSERVATION = NOT_CLAIMED
VIRTUAL_COSMOGENESIS = CONDITIONAL_MODEL
PHYSICAL_BIG_BANG_IDENTITY = NOT_CLAIMED
INDEPENDENT_REPRODUCTION = OPEN
GLOBAL_PASS = NOT_CLAIMED
FINAL_PASS = NOT_CLAIMED
EFFECT_ACK_DONE = NOT_CLAIMED
EFFECT_STATE = EFFECT_ACK_CONTINUE
\end{lstlisting}

\begin{thebibliography}{99}

\bibitem{cernsm}
CERN, \emph{The Standard Model}.\newline
\url{https://home.cern/science/physics/standard-model/}

\bibitem{cernhiggs}
CERN, \emph{The Higgs boson}.\newline
\url{https://home.cern/science/physics/higgs-boson/}

\bibitem{atlas2012}
ATLAS Collaboration, \emph{Observation of a new particle in the search for the
Standard Model Higgs boson with the ATLAS detector at the LHC}, Physics Letters
B 716 (2012). DOI: \href{https://doi.org/10.1016/j.physletb.2012.08.020}{10.1016/j.physletb.2012.08.020}.

\bibitem{cms2012}
CMS Collaboration, \emph{Observation of a new boson at a mass of 125 GeV with
the CMS experiment at the LHC}, Physics Letters B 716 (2012). DOI:
\href{https://doi.org/10.1016/j.physletb.2012.08.021}{10.1016/j.physletb.2012.08.021}.

\bibitem{ligo2016}
B. P. Abbott et al., \emph{Observation of Gravitational Waves from a Binary
Black Hole Merger}, Physical Review Letters 116, 061102 (2016). DOI:
\href{https://doi.org/10.1103/PhysRevLett.116.061102}{10.1103/PhysRevLett.116.061102}.

\bibitem{deser1970}
S. Deser, \emph{Self-Interaction and Gauge Invariance}, General Relativity and
Gravitation 1, 9--18 (1970). DOI:
\href{https://doi.org/10.1007/BF00759198}{10.1007/BF00759198}.

\bibitem{donoghue1994}
J. F. Donoghue, \emph{General relativity as an effective field theory: The
leading quantum corrections} (1994).\newline
\url{https://arxiv.org/abs/gr-qc/9405057}

\bibitem{hkm1976}
S. W. Hawking, A. R. King and P. J. McCarthy, \emph{A new topology for curved
space-time which incorporates the causal, differential, and conformal
structures}, Journal of Mathematical Physics 17 (1976). DOI:
\href{https://doi.org/10.1063/1.522874}{10.1063/1.522874}.

\bibitem{malament1977}
D. B. Malament, \emph{The class of continuous timelike curves determines the
topology of spacetime}, Journal of Mathematical Physics 18 (1977). DOI:
\href{https://doi.org/10.1063/1.523436}{10.1063/1.523436}.

\bibitem{bombelli1987}
L. Bombelli, J. Lee, D. Meyer and R. D. Sorkin, \emph{Space-time as a causal
set}, Physical Review Letters 59 (1987). DOI:
\href{https://doi.org/10.1103/PhysRevLett.59.521}{10.1103/PhysRevLett.59.521}.

\bibitem{oreshkov2012}
O. Oreshkov, F. Costa and C. Brukner, \emph{Quantum correlations with no causal
order}, Nature Communications 3, 1092 (2012). DOI:
\href{https://doi.org/10.1038/ncomms2076}{10.1038/ncomms2076}.

\bibitem{procopio2015}
L. M. Procopio et al., \emph{Experimental superposition of orders of quantum
gates}, Nature Communications 6, 7913 (2015). DOI:
\href{https://doi.org/10.1038/ncomms8913}{10.1038/ncomms8913}.

\end{thebibliography}

\vfill
\begin{center}
\small
Autor: Ingolf Lohmann · \QIKVRT{} VRTCore SMG H5 · 2. August 2026\\
\code{q.e.d.} gilt ausschließlich für die benannten formalen Modellsätze.
\end{center}

\end{document}
